% Standalone problem document, generated from the adjacent README.md.
% Compile with XeLaTeX. Mathematical content is maintained in Markdown.
\documentclass[11pt,a4paper]{article}
\usepackage[fontsize=10.5pt]{scrextend}
\usepackage[margin=22mm,headheight=15pt,headsep=9mm,footskip=11mm]{geometry}
\usepackage{fontspec}
\setmainfont{texgyrepagella}[Extension=.otf,UprightFont=*-regular,BoldFont=*-bold,ItalicFont=*-italic,BoldItalicFont=*-bolditalic]
\setsansfont{texgyreheros}[Extension=.otf,UprightFont=*-regular,BoldFont=*-bold,ItalicFont=*-italic,BoldItalicFont=*-bolditalic]
\setmonofont{texgyrecursor}[Extension=.otf,UprightFont=*-regular,BoldFont=*-bold,ItalicFont=*-italic,BoldItalicFont=*-bolditalic,Scale=MatchLowercase]
\usepackage{amsmath,amssymb}
\usepackage{unicode-math}
\setmathfont{texgyrepagella-math.otf}
\usepackage{xcolor,microtype,enumitem,needspace,fancyhdr}
\usepackage{bookmark}
\definecolor{ink}{HTML}{17344B}
\definecolor{accent}{HTML}{197D80}
\definecolor{muted}{HTML}{596773}
\hypersetup{colorlinks=true,linkcolor=accent,urlcolor=accent,pdftitle={AA-01 - Deciding
accurate evaluability of real
polynomials},pdfauthor={Open Problems in Numerical Linear Algebra}}
\urlstyle{same}
\setlength{\parindent}{0pt}
\setlength{\parskip}{5pt plus 1pt minus 1pt}
\linespread{1.04}
\setlength{\emergencystretch}{3em}
\widowpenalty=10000
\clubpenalty=10000
\displaywidowpenalty=10000
\allowdisplaybreaks[1]
\setlist{nosep,leftmargin=1.5em}
\providecommand{\tightlist}{\setlength{\itemsep}{0pt}\setlength{\parskip}{0pt}}
\providecommand{\pandocbounded}[1]{#1}
\pagestyle{fancy}
\fancyhf{}
\fancyhead[L]{\sffamily\footnotesize\color{muted}OPEN PROBLEMS IN NUMERICAL LINEAR ALGEBRA}
\fancyhead[R]{\sffamily\bfseries\footnotesize\color{accent}AA-01}
\fancyfoot[L]{\parbox[b]{0.9\textwidth}{\sffamily\fontsize{8}{10}\selectfont\color{muted}Editorial ratings; a bounded search cannot certify that a problem remains unsolved.\\Literature check: 2026-09-11}}
\fancyfoot[R]{\sffamily\footnotesize\color{muted}\thepage}
\renewcommand{\headrulewidth}{0.3pt}
\renewcommand{\headrule}{\hbox to\headwidth{\color{accent}\leaders\hrule height \headrulewidth\hfill}}
\makeatletter
\renewcommand{\section}{\Needspace{5\baselineskip}\@startsection{section}{1}{0pt}{12pt plus 2pt minus 2pt}{4pt}{\sffamily\bfseries\large\color{ink}}}
\renewcommand{\subsection}{\Needspace{4\baselineskip}\@startsection{subsection}{2}{0pt}{9pt plus 1pt minus 1pt}{3pt}{\sffamily\bfseries\normalsize\color{ink}}}
\makeatother
\setcounter{secnumdepth}{0}
\begin{document}
{\sffamily\small\color{muted}Arithmetic and complexity\par}
\vspace{3pt}
{\raggedright\hyphenpenalty=10000\exhyphenpenalty=10000\sffamily\bfseries\fontsize{21}{24}\selectfont\color{ink}Deciding
accurate evaluability of real polynomials\par}
\vspace{7pt}
{\sffamily\small\textbf{Difficulty:} extreme\quad\textbf{Importance:} broadly
interesting\par}
{\sffamily\small\bfseries\color{accent}Status: Solved\par}
\vspace{4pt}
{\color{accent}\hrule height 0.8pt}
\vspace{6pt}
\textbf{Rating rationale:} Extreme because the existence of any accurate
computation tree is a general decidability problem beyond checking a
supplied algorithm; broad importance includes numerical stability,
symbolic computation and automated algorithm design.

\subsection{Independently reviewed resolution -
2026-09-11}\label{independently-reviewed-resolution---2026-09-11}

Theorem 1.1 gives an always-halting decision procedure for the exact
constant-free finite-tree model in the statement. In every signed
ordering chart \(x_{\pi(i)}=\sigma_i(y_1+\cdots+y_i)\), the
coefficientwise absolute majorant of \(q(y)=p(x(y))\) must be bounded by
\(C|q(y)|\) on \(y\geq0\). This finite real-quantifier criterion is
necessary and sufficient; Sections 2-4 prove the equivalence and
construct an evaluator when it holds. Both independent reviews cover
comparisons, branching, stored-value reuse and arbitrary independent
rounding errors.

\textbf{Author:} Matthew J. Colbrook, Department of Applied Mathematics
and Theoretical Physics, University of Cambridge.
\href{https://github.com/ajt60gaibb/OpenProblemsInNLA/blob/main/references/colbrook-arithmetic-2026-09-11/manuscripts/AA-01.pdf}{Complete
manuscript},
\href{https://github.com/ajt60gaibb/OpenProblemsInNLA/blob/main/references/colbrook-arithmetic-2026-09-11/verification/reviews/AA-01-review.md}{independent
review}, and
\href{https://github.com/ajt60gaibb/OpenProblemsInNLA/blob/main/references/colbrook-arithmetic-2026-09-11/README.md}{submission
record}. AI assistance is disclosed. Agent verification is not external
human peer review or formal certification; no novelty or priority claim
is made.

A
\href{https://github.com/ajt60gaibb/OpenProblemsInNLA/blob/main/references/colbrook-arithmetic-2026-09-11/verification/reviews/AA-01-second-review.md}{second
independent proof review} also passes. The described experimental
programs were not retained; their run counts and compiler/solver claims
are not reproduced or certified here. The full mathematical proof,
rather than those logs, supports the resolution. The ratings and
rationale above are historical assessments of the original open target.

\subsection{Separate reviewed proof by George Stepaniants ---
2026-09-11}\label{separate-reviewed-proof-by-george-stepaniants-2026-09-11}

\textbf{Affirmative resolution by George Stepaniants}, Department of
Computing and Mathematical Sciences, California Institute of Technology.
The
\href{https://github.com/ajt60gaibb/OpenProblemsInNLA/blob/main/arithmetic-and-complexity/AA-01/solution.tex}{complete
manuscript}
(\href{https://github.com/ajt60gaibb/OpenProblemsInNLA/blob/main/arithmetic-and-complexity/AA-01/solution.pdf}{PDF}),
\textbf{Theorem 2.1 and Corollary 7.1}, gives an always-halting decision
algorithm for exactly the constant-free finite-tree model below.

For each signed ordering of the input magnitudes, write the inputs in
nonnegative gap coordinates and expand the transformed polynomial
\(q_T\). Accurate evaluability is equivalent to a uniform bound
\(W_T\le C_T|q_T|\), where \(W_T\) is its coefficientwise absolute
polynomial. The necessity proof couples actual executions, including
error-dependent branches and stored values. A termwise evaluator proves
sufficiency, including zeros and tied magnitudes. Finitely many
real-quantifier-elimination tests decide the criterion and construct an
evaluator on positive instances. No running-time bound or additional
arithmetic primitive is assumed.

A separate Codex agent checked the full proof against the original
target and returned \textbf{PASS} on 11 September 2026:
\href{https://github.com/ajt60gaibb/OpenProblemsInNLA/blob/main/references/stepaniants-2026-09-11/verification/reviews/AA-01-review.md}{independent
review}. The draft was developed in a ChatGPT conversation. Verification
is independent automated-agent review, not external human peer review or
formal proof certification. The original statement and earlier
literature checks are retained; the ratings above are historical.
\href{https://github.com/ajt60gaibb/OpenProblemsInNLA/blob/main/references/stepaniants-2026-09-11/README.md}{Authorship,
source record and submission details}.

\textbf{Integration note --- 2026-09-11.} This page retains both
separately attributed, independently reviewed proofs. Stepaniants's
\href{https://github.com/ajt60gaibb/OpenProblemsInNLA/pull/68}{existing
submission PR \#68} was submitted while this exact target was open;
Colbrook's
\href{https://github.com/ajt60gaibb/OpenProblemsInNLA/pull/89}{PR \#89}
has since been accepted into main. Both manuscripts and their review
records are preserved, as requested by the maintainer. The mathematical
statement below is unchanged.

\subsection{Problem statement}\label{problem-statement}

Consider finite arithmetic computation trees with exact real inputs
\(x\in\mathbb R^n\). Arithmetic nodes use binary addition, subtraction,
or multiplication, each returning \((a\mathbin{\circ}b)(1+\delta_j)\).
Negation and comparisons (\(<,=,>\)) of available quantities are exact,
and branching is allowed. There are no constant input nodes. Each
operation may have a different, arbitrary error \(\delta_j\); equal
operands need not produce equal errors. Previously computed values may
be stored and reused. The tree describes control flow, not a formula
that recomputes each occurrence of a subexpression.

For \(p\in\mathbb Z[x_1,\ldots,x_n]\) with \(p(0)=0\), call such a tree
\(P\) accurate if \(P(x,0)=p(x)\) and

\[
\forall\eta\in(0,1)\ \exists u\in(0,1)\quad
\forall x\in\mathbb R^n\ \forall\delta\in[-u,u]^{N(P)}:
\quad |P(x,\delta)-p(x)|\le\eta|p(x)|,
\]

where \(N(P)\) numbers all rounded nodes; errors on unused branches are
ignored. The same tree must work for every \(\eta\), and \(u\) may
depend on \(p,P,\eta\) but not on \(x\). The condition includes exact
output at zeros of \(p\).

\subsubsection{Question}\label{question}

Is there a Turing algorithm which, from the finite integer coefficient
list of any such \(p\), always halts and decides whether an accurate
tree exists? No bound on the decision algorithm's running time is
prescribed.

This is a basic obstruction to automatically constructing accurate
algorithms for determinants and other polynomial expressions associated
with structured matrices. Checking a supplied fixed tree is already
decidable; deciding whether any suitable tree exists is the question.

\subsection{References and status}\label{references-and-status}

J. Demmel, I. Dumitriu, O. Holtz, and P. Koev,
\href{https://people.eecs.berkeley.edu/~demmel/Demmel_pubs_07_11_final/B15_ActaNumerica08.pdf}{\emph{Accurate
and Efficient Expression Evaluation and Linear Algebra}}, \emph{Acta
Numerica} 17 (2008), 87--145, Definitions 3.3--3.4, §3.3 and especially
§3.3.7. The finite-tree, integer-coefficient formulation fixes the
traditional real-arithmetic version of their decision question. Theorem
3.12 answers the corresponding whole-space complex question, not the
real one. The original Demmel--Dumitriu--Holtz paper,
\href{https://arxiv.org/pdf/math/0508350v2}{\emph{Toward accurate
polynomial evaluation in rounded arithmetic}, v2}, §§2.2--2.7,
explicitly specifies uniformly bounded running time, exact comparisons,
and separate rounding errors, clarifying the computation-tree
conventions. Demmel's
\href{https://simons.berkeley.edu/events/when-accurate-efficient-expression-evaluation-linear-algebra-possible}{21
October 2025 Simons seminar abstract} still identifies the broader
accurate-evaluation decision procedure as open. Searches on 2026-09-08
for accurate polynomial evaluability, real-arithmetic decidability, and
follow-ups to the cited authors found no complete decision procedure or
undecidability theorem for the displayed model. The recent talk does not
separately specify every tree convention used here.

\subsection{Earlier status check ---
2026-09-10}\label{earlier-status-check-2026-09-10}

Rechecked
\href{https://people.eecs.berkeley.edu/~demmel/Demmel_pubs_07_11_final/B15_ActaNumerica08.pdf}{Demmel
et al., §3.3.7} and
\href{https://simons.berkeley.edu/events/when-accurate-efficient-expression-evaluation-linear-algebra-possible}{Demmel's
October 2025 Simons abstract}, and searched for later real-polynomial
decision procedures. The broad real accurate-evaluation question remains
explicitly open in the seminar abstract; no complete algorithm or
undecidability theorem for the displayed finite-tree model was located.
The complex-domain characterization and fixed-tree verification do not
answer this existence question.
\end{document}
